\section{Resultados}
Los resultados de prueba corresponden a 2v1, defensor a potencia 50 y arquero
activo. Se distingue el éxito de cada habilidad de la tasa de gol del ataque
completo; las curvas de entrenamiento no sustituyen la evaluación con pesos fijos.

\subsection{Aprendizaje desde cero y escalera de algoritmos}
\begin{figure*}[t]
    \centering
    \includegraphics[width=\textwidth]{pilares/pilar3_curvas.png}
    \caption{Retorno base de persecución: media y banda $\pm1\sigma$ entre tres
    semillas, con ventana móvil de 50 episodios. Todas estas variantes entrenan con PBRS.}
    \label{fig:curvas}
\end{figure*}

La Fig.~\ref{fig:curvas} compara la evolución de los métodos de valor y
actor-crítico. DQN y Double DQN alcanzan pronto retornos elevados; PPO necesita
más interacciones y A2C mantiene retornos inferiores en el presupuesto estudiado.
La inicialización por BC comienza cerca del nivel final, sin incluir en el eje
el costo de obtener demostraciones y preentrenar.

\begin{table}[H]
    \centering
    \caption{Persecución: retorno final $\bar G_0\pm\sigma$ entre semillas,
    pasos $t_{80}$ y éxito en prueba. Cada agente profundo reúne 600 episodios;
    no hay evaluación tabular equivalente.}
    \label{tab:escalera}
    \footnotesize
    \setlength{\tabcolsep}{3pt}
    \begin{tabular}{@{}lrrr@{}}
        \toprule
        Método & $\bar G_0\pm\sigma$ & $t_{80}$ & Éxito (\%) \\
        \midrule
        Q tabular & $8{,}55\pm0{,}08$ & 10,0k & --- \\
        DQN & $8{,}14\pm0{,}34$ & 9,0k & 99,8 \\
        Double DQN & $8{,}26\pm0{,}21$ & 10,3k & 97,0 \\
        A2C & $2{,}25\pm0{,}80$ & --- & 72,0 \\
        PPO & $7{,}92\pm0{,}58$ & 25,0k & 98,7 \\
        BC humano + PPO & $8{,}99\pm0{,}01$ & 0 & 99,8 \\
        \midrule
        PPO sin shaping & $3{,}92\pm1{,}71$ & --- & 64,2 \\
        PPO shaping ingenuo & $2{,}47\pm1{,}20$ & --- & 40,5 \\
        \bottomrule
    \end{tabular}
\end{table}

La Tabla~\ref{tab:escalera} resume el último 20\,\% de episodios de cada corrida.
El umbral común es $G_{80}=7{,}02$; ``---'' en $t_{80}$ indica que ninguna semilla
lo alcanza. DQN requiere 9k pasos frente a 25k de PPO. BC + PPO presenta la menor
dispersión final. El baseline tabular tiene observaciones y acciones reducidas,
por lo que estas cifras no aíslan exclusivamente el efecto del aproximador.

\begin{figure*}[t]
    \centering
    \includegraphics[width=0.9\textwidth]{pilares/pilar3_perdidas.png}
    \caption{Pérdida TD, valor Q medio, pérdida del crítico y entropía durante
    el entrenamiento: media y banda $\pm1\sigma$ entre tres semillas.}
    \label{fig:perdidas}
\end{figure*}

La Fig.~\ref{fig:perdidas} complementa los retornos con la evolución de la
optimización. Las pérdidas TD y del crítico corresponden a objetivos distintos
y no permiten ordenar directamente las familias. Asimismo, una diferencia de
valores Q entre DQN y Double DQN no demuestra por sí sola sobreestimación sin
compararlos con los retornos de referencia.

\subsection{PBRS y shaping ingenuo}
\begin{figure*}[t]
    \centering
    \includegraphics[width=\textwidth]{pilares/pilar2_shaping.png}
    \caption{Ablación del shaping en tres semillas: retorno base en persecución,
    éxito móvil en conducción y suma no descontada de recompensas recibidas.
    Las variantes de conducción parten del mismo modelo BC.}
    \label{fig:shaping}
\end{figure*}

La Fig.~\ref{fig:shaping} separa la recompensa recibida del cumplimiento de la
tarea. En persecución, PBRS obtiene 98,7\,\% de éxito, frente a 64,2\,\% sin
shaping y 40,5\,\% con el ingenuo. En conducción, las tasas son 99,2\,\%,
98,0\,\% y 92,7\,\%, respectivamente, con 600 episodios por variante. El shaping
ingenuo aumenta la recompensa acumulada durante el entrenamiento y reduce el
éxito de prueba: la señal favorece comportamientos que no mejoran la tarea.
La identidad telescópica se aplica al retorno descontado, no a la suma sin
descuento del panel derecho.

\subsection{Imitación y ajuste fino por habilidad}
\begin{table}[H]
    \centering
    \caption{Éxito en 200 episodios por modelo (\%). En definición, éxito significa gol.}
    \label{tab:habilidades}
    \footnotesize
    \setlength{\tabcolsep}{3pt}
    \begin{tabular}{@{}lrrr@{}}
        \toprule
        Variante & Persec. & Conduc. & Defin. \\
        \midrule
        BC humano & 100,0 & 13,0 & 57,5 \\
        DAgger humano & 100,0 & 53,5 & 58,0 \\
        BC experto programado & --- & 99,5 & --- \\
        DAgger experto programado & --- & 100,0 & --- \\
        Imitación + DQN & 100,0 & 70,5 & 83,0 \\
        Imitación + A2C & 100,0 & 100,0 & 52,0 \\
        Imitación + PPO & 100,0 & 98,0 & 63,5 \\
        PPO desde zona de tiro & --- & --- & 86,0 \\
        \bottomrule
    \end{tabular}
\end{table}

La Tabla~\ref{tab:habilidades} recoge las comparaciones de las diapositivas.
DAgger humano aumenta conducción en 40,5 puntos porcentuales respecto a BC,
pero apenas cambia definición. El experto programado alcanza al menos 99,5\,\%
en conducción. Tras imitación, los tres algoritmos capturan el balón en todos
los episodios de prueba; en conducción destacan A2C/PPO y en definición DQN.
PPO entrenado desde la zona de tiro obtiene 86,0\,\% en la habilidad aislada,
lo que no equivale a esa tasa de gol dentro de la jerarquía.

\begin{figure*}[t]
    \centering
    \includegraphics[width=\textwidth]{imitacion/aprendizaje_subtareas.png}
    \caption{Comparación por subpolítica: retorno base desde cero (izquierda),
    retorno base tras imitación y ajuste fino (centro), y éxito del modelo
    ajustado en 200 episodios de prueba con IC de Wilson al 95\,\% (derecha).
    Desde cero: tres semillas en persecución y dos en las otras tareas, con
    bandas $\pm1\sigma$ donde hay réplicas y ventana de 50 episodios.
    Ajuste fino: una corrida por algoritmo,
    con ventana de 100 episodios en persecución/conducción y 200 en definición.}
    \label{fig:ajuste_subtareas}
\end{figure*}

La Fig.~\ref{fig:ajuste_subtareas} reúne las comparaciones de retorno y éxito
por habilidad. Las dos primeras columnas usan la misma recompensa base y escala
vertical; cambia la inicialización, no la definición del retorno. Tras imitación,
A2C y PPO mantienen retornos altos en conducción, mientras DQN cae y se recupera
parcialmente. En definición destaca DQN, con 83,0\,\% de gol en prueba. Las series
desde cero y de ajuste tienen presupuestos y condiciones de entrenamiento
distintos: muestran su evolución, pero no aíslan el efecto causal de la imitación
ni acreditan estabilidad entre semillas del ajuste fino.

\subsection{Ataque jerárquico completo}
\begin{figure}[t]
    \centering
    \includegraphics[width=\columnwidth]{top_policies/jerarquia_final.png}
    \caption{Goles del ataque 2v1 en 200 episodios por variante; barras de error:
    IC de Wilson al 95\,\%. La línea discontinua representa la meta del 80\,\%.}
    \label{fig:jerarquia}
\end{figure}

La Fig.~\ref{fig:jerarquia} actualiza la comparación con las variantes recientes
mostradas en la exposición. El coordinador DAgger humano obtiene 30,5\,\%; el
ajuste del sistema completo alcanza 26,5\,\% con PPO, 16,0\,\% con A2C y 1,0\,\%
con DQN. La política plana no convierte goles en esta prueba.

Con habilidades especializadas, BC del coordinador obtiene 41,0\,\%, DAgger
44,5\,\% y el ajuste posterior con PPO 47,0\,\%. El currículo de definición
alcanza 51,5\,\%. La variante de traspaso con ventaja obtiene 45,0\,\%; al
entrenar definición dentro de esa misma jerarquía sube a 65,0\,\%, equivalente
a $130/200$ goles e IC 95\,\% de 58,2--71,3\,\%. Es el mejor resultado registrado
en esta comparación, pero sigue por debajo del objetivo del 80\,\%. Estas
variantes tienen presupuestos y etapas de entrenamiento distintos: la comparación
no atribuye toda la mejora a un único cambio.

\begin{figure}[t]
    \centering
    \includegraphics[width=\columnwidth]{04_jerarquico/reward_y_goles.png}
    \caption{Retorno base durante el entrenamiento de definición dentro de
    la jerarquía, con media móvil de 1\,000 episodios.}
    \label{fig:reward_jerarquico}
\end{figure}

La Fig.~\ref{fig:reward_jerarquico} muestra un aumento del retorno durante el
entrenamiento, con fluctuaciones y una estabilización final alrededor de
$G_0=1{,}1$. La tendencia indica una mejora de la recompensa base, aunque
persisten retrocesos durante el aprendizaje.
